\documentclass{standalone}
\usepackage{tikz}
\usetikzlibrary{shapes.geometric, arrows, positioning}

\begin{document}

\tikzstyle{iniciofim} = [ellipse, draw, fill=gray!20, text width=3cm, text centered, minimum height=1cm]
\tikzstyle{processo} = [rectangle, draw, fill=orange!15, text width=4.2cm, text centered, rounded corners, minimum height=1cm]
\tikzstyle{decisao} = [diamond, draw, fill=green!10, text width=3.5cm, text centered, aspect=2]
\tikzstyle{seta} = [->, thick]

\begin{tikzpicture}[node distance=1.6cm]

\node (inicio) [iniciofim] {Início};
\node (leitura) [processo, below of=inicio] {Importar dados do SD para o Python};
\node (estatistica) [processo, below of=leitura] {Calcular média, mediana, desvio padrão e IQR};
\node (outliers) [processo, below of=estatistica] {Detectar e remover outliers};
\node (correlacao) [processo, below of=outliers] {Calcular correlação de Pearson};
\node (graficos) [processo, below of=correlacao] {Gerar gráficos e relatórios};
\node (fim) [iniciofim, below of=graficos] {Fim da análise};

\draw [seta] (inicio) -- (leitura);
\draw [seta] (leitura) -- (estatistica);
\draw [seta] (estatistica) -- (outliers);
\draw [seta] (outliers) -- (correlacao);
\draw [seta] (correlacao) -- (graficos);
\draw [seta] (graficos) -- (fim);

\end{tikzpicture}

\end{document}

